% Part: normal-modal-logic
% Chapter: axioms-systems
% Section: normal-logics

\documentclass[../../../include/open-logic-section]{subfiles}

\begin{document}

\olfileid{nml}{prf}{nor}

\olsection{المنطقات الجهية النظامية}

ليس من السهل توصيف كل مجموعة من !!{formula}s الجهية بأنها مجموعة
!!{formula}s القابلة للاشتقاق من مجموعة بديهيات. فنحن نريد للمنطقات
الجهية أن تكون حسنة السلوك. وأولا، ينبغي بالطبع أن يظل كل ما نستطيع
أن !!{derive} في منطق القضايا الكلاسيكي !!{derivable}، مع مراعاة أن
!!{formula}s قد تحتوي الآن أيضا على
\iftag{prvBox}{$\Box$\iftag{prvDiamond}{
    و~}{}}{}\iftag{prvDiamond}{$\Diamond$}{}. ولذلك نشترط أن يشتمل
المنطق الجهوي على جميع أمثلة القضايا الصادقة صدقا تحليليا، وأن يكون
مغلقا تحت القياس الاستثنائي.

\begin{defn}
  \emph{المنطق الجهوي} هو مجموعة~$\Sigma$ من !!{formula}s الجهية بحيث
  \begin{enumerate}
  \item تشتمل على جميع القضايا الصادقة صدقا تحليليا، و
  \item تكون مغلقة تحت الإحلال، أي إذا كان $!A \in \Sigma$ وكانت
    $!D_1$، \dots، $!D_n$ هي !!{formula}s، فعندئذ
    \[
    \SSubst{!A}{\subst{!D_1}{p_1}, \dots, \subst{!D_n}{p_n}} \in \Sigma,
    \]
  \item تكون مغلقة تحت \emph{القياس الاستثنائي}، أي إذا كان $!A$
    و$!A \lif !B \in \Sigma$، فعندئذ $!B \in \Sigma$.
  \end{enumerate}
\end{defn}

لكي نستعمل الدلالة العلاقية للمنطقات الجهية، علينا أيضا أن نشترط
اشتمالها على جميع !!{formula}s الصالحة في جميع النماذج الجهية. ويتبين
أن هذا الشرط يتحقق متى كانت جميع أمثلة \Ax{K}
\iftag{prvDiamond}{و~\Dual{}}{} !!{derivable}، وكان كلما كانت
!!a{formula}~$!A$ !!{derivable} كانت~$\Box !A$ كذلك. ويسمى المنطق
الجهوي الذي يحقق هذه الشروط \emph{نظاميا}. (وتوجد بالطبع منطقات
جهية غير نظامية، لكن النماذج العلاقية المعتادة لا تلائمها.)

\begin{defn}
  يكون المنطق الجهوي $\Sigma$ \emph{نظاميا} إذا اشتمل على
  \begin{align*}
    \tag{\Ax{K}} & \Box(p \lif q) \lif (\Box p \lif \Box q),
    \iftag{prvDiamond}{\\
      \tag{\Dual} & \Diamond p \liff \lnot\Box\lnot p}{}
  \end{align*}
  وكان مغلقا تحت \emph{الإلزام}، أي إذا كان $!A \in
  \Sigma$، فعندئذ $\Box !A \in \Sigma$.
\end{defn}

لاحظ أن اللزوم الصادق صدقا تحليليا «دقيق» بما يكفي ليحفظ
\emph{الصدق في عالم}، بينما لا تحفظ القاعدة \Nec{} إلا \emph{الصدق
في نموذج} (ومن ثم تحفظ أيضا الصلاحية في إطار أو في فئة أطر).

\begin{prop}\ollabel{prop:rk}
  كل منطق جهوي نظامي مغلق تحت القاعدة \RK،
  \begin{prooftree}
    \AxiomC{$!A_1 \lif (!A_2 \lif \cdots (!A_{n-1} \lif !A_n)\cdots)$}
    \RightLabel{\RK}
    \UnaryInfC{$\Box!A_1 \lif (\Box!A_2 \lif \cdots (\Box!A_{n-1}
      \lif \Box!A_n)\cdots).$}
  \end{prooftree}
\end{prop}

\begin{proof}
  بالاستقراء على~$n$: إذا كان $n = 1$، فالقاعدة ليست سوى \Nec، وكل
  منطق جهوي نظامي مغلق تحت \Nec.

  ولنفترض الآن أن النتيجة صحيحة لـ$n-1$؛ ونبين صحتها لـ$n$.

  افترض أن
  \begin{align*}
  & !A_1 \lif (!A_2 \lif \cdots (!A_{n-1} \lif !A_n)\cdots) \in \Sigma
  \intertext{وبفرض الاستقراء، لدينا}
  & \Box!A_1 \lif (\Box!A_2 \lif \cdots \Box(!A_{n-1} \lif !A_n)\cdots)
  \in \Sigma
  \intertext{ولأن $\Sigma$ منطق جهوي نظامي، فإنه يشتمل على جميع أمثلة
    $\Ax{K}$، ومنها بوجه خاص}
  & \Box(!A_{n-1} \lif !A_n) \lif (\Box!A_{n-1} \lif \Box!A_n) \in \Sigma
  \intertext{وباستخدام القياس الاستثنائي وأمثلة مناسبة لقضايا صادقة صدقا تحليليا نحصل على}
  & \Box!A_1 \lif (\Box!A_2 \lif \cdots (\Box!A_{n-1}
  \lif \Box!A_n)\cdots) \in \Sigma.
  \end{align*}
\end{proof}

\begin{prop}\ollabel{prop:notDiamondBot}
  يشتمل كل منطق جهوي نظامي $\Sigma$ على~$\lnot\Diamond\lfalse$.
\end{prop}

\begin{prob}
  أثبت \olref[nml][prf][nor]{prop:notDiamondBot}.
\end{prob}

\begin{prop}
  لتكن $!A_1$، \dots، $!A_n$ هي !!{formula}s. عندئذ يوجد أصغر
  منطق جهوي نظامي $\Sigma$ يشتمل على جميع أمثلة $!A_1$، \dots، $!A_n$.
\end{prop}

\begin{proof}
  إذا أعطيت $!A_1$، \dots، $!A_n$، فعرّف $\Sigma$ بأنه تقاطع جميع
  المنطقات الجهية النظامية التي تشتمل على جميع أمثلة $!A_1$، \dots،
  $!A_n$. وهذا التقاطع غير خال، لأن $\Frm[L]$، وهي مجموعة جميع
  !!{formula}s، منطق جهوي من هذا القبيل.
\end{proof}

\begin{defn}
يسمى أصغر منطق جهوي نظامي يشتمل على $!A_1$، \dots، $!A_n$
\emph{نسقا جهيا}، ويرمز إليه بـ$\Log{K} !A_1 \dots
!A_n$. ويرمز إلى أصغر منطق جهوي نظامي بـ~\Log{K}.
\end{defn}

\end{document}
